\uuid{6B3h}
\exo7id{4024}
\titre{exo7 4024}
\auteur{quercia}
\organisation{exo7}
\datecreate{2010-03-11}
\isIndication{false}
\isCorrection{true}
\chapitre{Développement limité}
\sousChapitre{Applications}
\cpge{mpsi-mp2i.analyse-asymptotique.b, pcsi.analyse-asymptotique.b, ptsi.analyse-asymptotique.b, mpsi-mp2i.limites-continuite.a, pcsi.limites-continuite.a, ptsi.limites-continuite.a, bcpst1.developpements-limites.a, bcpst1.limites-continuite.a}
\module{Analyse}
\niveau{L1}
\difficulte{}

\contenu{
\question{
Soient $a_1,\dots,a_n \in \R$.
CNS pour que $\sum_{k=1}^n a_k \mathrm{cotan}(kx)$ ait une limite finie en 0 ?
}
\reponse{
$\sum_{k = 1}^n \frac {a_k}k = 0 \Rightarrow  L = \frac 13 \sum_{k = 1}^n a_k$.
}
}
